%! Simplifying Rational Expressions — Unit 1, Lesson 1.7 — Slides (Beamer)
\documentclass[aspectratio=169, 11pt]{beamer}
\usepackage{atda-beamer}   % provides \royalheader{} and \sectionlabel[color]{}

\begin{document}

% TITLE SLIDE (CourseName is not defined in beamer, so the name is literal)
{
\setbeamercolor{background canvas}{bg=royal}
\begin{frame}[plain]
  \vspace{2.8cm}\hspace{0.55cm}%
  \begin{minipage}{0.9\textwidth}
    {\color{goldacc}\bfseries\small LESSON 1.7}\\[0.5cm]
    {\color{white}\bfseries\LARGE Simplifying Rational Expressions}\\[1.2cm]
    {\color{mistmid}\small Algebra, Trigonometry, and Data Analysis~$\cdot$~Unit 1}
  \end{minipage}
\end{frame}
}

% ── HOOK ──────────────────────────────────────────────────────────────────────
\begin{frame}[t, plain]
  \royalheader{The Booster Club's T-Shirt Order}
  \sectionlabel[goldacc]{hook --- does a shirt really cost \$246?}
  A \$$240$ setup fee plus \$$6$ a shirt, so $x$ shirts cost $6x+240$ dollars:
  \[ A(x) = \frac{6x+240}{x} \quad \text{dollars per shirt.} \]
  \begin{block}{A student cancels the $x$'s and gets \$246}
    Test it. $A(10)=\$30$, \quad $A(40)=\$12$, \quad $A(240)=\$7$. \\
    Not \$246 --- not even close. And what happens at $x=0$?
  \end{block}
\end{frame}

% ── THE WARM-UP IS THE LESSON ─────────────────────────────────────────────────
\begin{frame}[t, plain]
  \royalheader{Plus Signs Do Not Cancel}
  \sectionlabel{warm-up item 3, in numbers you already trust}
  \[ \frac{12}{18} = \frac{6 \cdot 2}{6 \cdot 3} = \frac{2}{3}
     \qquad\qquad \frac{3 \cdot 4}{3 \cdot 5} = \frac{4}{5} \]
  \[ \frac{3+4}{3+5} = \frac{7}{8}
     \;{\color{goldacc}\ne}\; \frac{4}{5} \]
  \begin{block}{The rule for the whole lesson}
    You may divide out only common \textbf{factors} --- things being \emph{multiplied}. Terms
    joined by $+$ or $-$ never cancel.
  \end{block}
\end{frame}

% ── EXCLUDED VALUES ───────────────────────────────────────────────────────────
\begin{frame}[t, plain]
  \royalheader{What Breaks a Rational Expression}
  \sectionlabel{\texttt{A2.EO.1b} --- excluded values are Lesson 1.6}
  Set the \textbf{denominator} equal to zero, factor, and use the zero product property.
  \vspace{0.15cm}
  \renewcommand{\arraystretch}{1.6}
  \begin{tabularx}{\textwidth}{@{}p{3.0cm} p{5.4cm} X@{}}
    \textbf{Expression} & \textbf{Denominator $=0$} & \textbf{Excluded values} \\
    \hline
    $\frac{5}{x-3}$        & $x-3=0$              & $x=3$ \\
    $\frac{x+1}{x^2-4}$    & $(x+2)(x-2)=0$       & $x=-2$, \ $x=2$ \\
    $\frac{3x}{x^2+9}$     & no real solution     & {\color{goldacc}none at all} \\
  \end{tabularx}
  \vspace{0.25cm}

  Row 3 is Lesson 1.5's prime sum of squares, paying a dividend.
\end{frame}

% ── SIMPLIFYING ───────────────────────────────────────────────────────────────
\begin{frame}[t, plain]
  \royalheader{Simplifying --- Factor, Then Divide Out}
  \sectionlabel{factor the top, factor the bottom, cancel what they share}
  \[
  \begin{aligned}
    \frac{x^2-9}{x^2+7x+12}
      &= \frac{(x+3)(x-3)}{(x+3)(x+4)} \\[0.15cm]
      &= \frac{x-3}{x+4} \qquad {\color{goldacc}(x \ne -3, \; x \ne -4)}
  \end{aligned}
  \]
  \begin{block}{Simplest form}
    The top and bottom share no common factor except $1$. ``Simplify'' does not mean ``make
    shorter'' --- sometimes there is nothing to divide out.
  \end{block}
\end{frame}

% ── THE SCAR ──────────────────────────────────────────────────────────────────
\begin{frame}[t, plain]
  \royalheader{A Cancelled Factor Leaves a Scar}
  \sectionlabel[goldacc]{read the excluded values from the \emph{original} denominator}
  \[ \frac{x^2-9}{x^2+7x+12} \quad\text{at } x=-3 \;\text{ reads }\; \frac{0}{0}
     \quad\text{--- undefined} \]
  \[ \frac{x-3}{x+4} \quad\text{at } x=-3 \;\text{ happily returns }\; -6 \]
  \begin{block}{So they are equivalent --- but only where the original was defined}
    Cancelling $(x+3)$ hides the restriction; it does not remove it. Carry $x \ne -3$ along with
    the answer.
  \end{block}
\end{frame}

% ── THE ERROR OF THE DAY ──────────────────────────────────────────────────────
\begin{frame}[t, plain]
  \royalheader{The Error of the Day}
  \sectionlabel[goldacc]{third appearance in this unit}
  \[ \frac{x+3}{x+5} \;{\color{goldacc}\ne}\; \frac{3}{5}
     \qquad\text{at } x=1: \;\; \frac{4}{6} = \frac{2}{3} \]
  \begin{block}{The same misconception, three lessons running}
    Lesson 1.1: $(2m+3)^2 \ne 4m^2+9$. \quad Lesson 1.2: $\sqrt{9+16} \ne 3+4$. \quad
    Today: you cannot cancel across a $+$. \\[0.2cm]
    An operation does not reach inside a sum. \textbf{Refute it with a number, every time.}
  \end{block}
\end{frame}

% ── MULTIPLY AND DIVIDE ───────────────────────────────────────────────────────
\begin{frame}[t, plain]
  \royalheader{Multiplying and Dividing}
  \sectionlabel{\texttt{A2.EO.1a} --- factor first, cancel across}
  \[
  \begin{aligned}
    \frac{x+2}{x-5} \cdot \frac{x^2-25}{x^2-4}
      &= \frac{x+2}{x-5} \cdot \frac{(x+5)(x-5)}{(x+2)(x-2)}
       = \frac{x+5}{x-2}
  \end{aligned}
  \]
  \begin{block}{To divide, multiply by the reciprocal}
    \[ \frac{x^2-1}{3x} \div \frac{x+1}{6x^2}
       = \frac{(x+1)(x-1)}{3x} \cdot \frac{6x^2}{x+1} = 2x(x-1) \]
    Excluded values come from \textbf{every} denominator you ever wrote --- including the one you
    flipped, so $x \ne -1$ survives.
  \end{block}
\end{frame}

% ── ADD AND SUBTRACT ──────────────────────────────────────────────────────────
\begin{frame}[t, plain]
  \royalheader{Adding and Subtracting}
  \sectionlabel{same bottom? combine the tops --- then simplify}
  \[
  \begin{aligned}
    \frac{x^2}{x-3} - \frac{9}{x-3}
      &= \frac{x^2-9}{x-3} = \frac{(x+3)(x-3)}{x-3} = x+3 \qquad (x \ne 3)
  \end{aligned}
  \]
  \begin{block}{Different bottoms? Build a common one}
    \[ \frac{4}{x-1} + \frac{3}{x+2}
       = \frac{4(x+2) + 3(x-1)}{(x-1)(x+2)} = \frac{7x+5}{(x-1)(x+2)} \]
    Check at $x=2$: \; $\tfrac41 + \tfrac34 = \tfrac{19}{4}$, and $\tfrac{19}{4}$ comes back out.
  \end{block}
\end{frame}

% ── ACTIVITY LAUNCH ───────────────────────────────────────────────────────────
\begin{frame}[t, plain]
  \royalheader{Group Activity: The Drama Club Backdrop}
  \sectionlabel{groups of three --- everyone starts at Tier R}
  The panel's area and width are on the plan; nobody wrote down the length.
  \[ A(x) = x^2+3x-28 \;\text{ sq ft}, \qquad \text{width } = (x-4) \text{ ft.} \]
  \begin{block}{Same move as Lessons 1.3--1.5, one level up}
    Factor, then divide out. Then: the paint bill is $C(x)=2x^2+6x-56$ --- what is the cost
    \emph{per square foot}, and why is the answer a plain number?
  \end{block}
\end{frame}

% ── CLOSE ─────────────────────────────────────────────────────────────────────
\begin{frame}[t, plain]
  \royalheader{Exit Ticket}
  \sectionlabel[goldacc]{notes closed --- 5 minutes}
  \begin{enumerate}
    \setlength{\itemsep}{0.35cm}
    \item Simplify, with excluded values: \; $\dfrac{x^2-25}{x^2+8x+15}$.
    \item Multiply and simplify: \; $\dfrac{x^2-9}{x+2} \cdot \dfrac{x^2-4}{x-3}$.
    \item A student says $\dfrac{x+6}{x+2} = \dfrac62 = 3$. Test it at $x=2$.
  \end{enumerate}
  \vspace{0.3cm}
  {\color{royal}\bfseries That closes Unit 1 --- next up, the unit test.}
\end{frame}

\end{document}
